%\title{My two column CV}
%
% tccv (two columns curriculum vitae) is a LaTeX class inspired by
% the template found at latextemplates.com by Alessandro Plasmati.
%
% Create by Nicola Fontana, the original files can be downloaded from:
% http://dev.entidi.com/p/tccv/
%
% Modified *heavily* by Pooria Azimi:
% https://github.com/pooriaazimi/CV/
%
\documentclass{tccv}

\usepackage[english]{babel}
\usepackage{ulem}
\usepackage[object=vectorian]{pgfornament}
\usepackage{setspace}
\usepackage{multirow}
\usepackage{hhline}
\usepackage{fancyhdr}
\usepackage{array}
\usepackage{graphicx}
\usepackage{url}
\usepackage{layout}
\usepackage{svg}
\usepackage{epstopdf}

\fancyhf{}
\fancyfoot[L]{Pooria Azimi -- Curriculum Vit\ae}
\fancyfoot[C]{(\thepage)}
\fancyfoot[R]{\href{http://bit.ly/pooria-azimi-cv}{http://bit.ly/pooria-azimi-cv}}
\renewcommand\headrulewidth{0pt}
\renewcommand\footrulewidth{1pt}
\setlength{\footskip}{22pt}

\renewcommand\rmdefault{lmr}
\renewcommand\sfdefault{lmss}
\renewcommand\ttdefault{lmtt}

\pagestyle{fancy}

\renewcommand{\tilde}{\raise.17ex\hbox{$\scriptstyle\sim$}}
\newcommand{\plus}{\raisebox{.1\height}{\scalebox{.8}{+}}}


\begin{document}

\thispagestyle{empty}
\part{Pooria Azimi}





\section{Research Interests}


\begin{research_interest}

\item{Database and Information Retrieval\\[4pt]Systems}
     {}

%\item{Database and IR Systems}
%     {}

\item{Operating Systems}
     {File systems -- Microkernels}

\item{Concurrent Programming}
     {the ``Actor Model''}

\item{Human-Computer Interaction}
     {}

\end{research_interest}





%\vspace{6.5pt}



\section{Education}

{\bf \href{https://en.wikipedia.org/wiki/Amirkabir_University_of_Technology}{Amirkabir University of Technology}} (Tehran, Iran)

\hspace{7pt} B.Sc. in {\bf Computer Science}, 2014 {\it(expected)}

\hspace{7pt} GPA (last 60 CS credits): {\bf 16.5}/20


%\vspace{10pt}




\section{Awards and Honors}

\begin{list}{\labelitemi}{\leftmargin=1em}

\item
{\bf 2\textsuperscript{nd} place} at \href{http://www.data-mining-cup.de/en/review/dmc-2012/}{\mbox{13\textsuperscript{th} International Data Mining Cup} (Berlin)}
\\
{\it Our team created a bidding agent in Java for the \mbox{``online''} task, and used a combination of statistical models, \mbox{neural} networks and SVMs (with \href{http://www.mathworks.com/products/matlab/}{MATLAB} and \href{http://www.cs.waikato.ac.nz/ml/weka/}{Weka}) for predicting the results of the \mbox{``offline''} task.}



\item
Ranked 590 (among \tilde400,000 participants) in 2008 \mbox{National} Matriculation Exam ({\bf top \%0.2})


\end{list}











\section{Communication skills}


\begin{factlist}

\item{English}{Fluent}

\item{Persian}{Native speaker}

\item{German}{Basic understanding}

\end{factlist}


\vspace{-4pt}

\sectionline

\vspace{-10pt}


{\renewcommand{\arraystretch}{1.6}
\begin{table}[ph]
  \centering
	\begin{tabular}{C{3.5cm}|C{4cm}}
 	   {\sc TOEFL }i{\sc BT} & {\sc GRE General (\uline{expected}\textsuperscript{$*$})}\\[5pt]
	   \hline
        \vspace{-12pt}
 	   	\begin{tabular}{R{1.7cm}|L{0.9cm}}
			Reading & 30 \\
			Listening & 29\\
			Speaking & 22\textsuperscript{$**$} \\
			Writing & 28\\
			\hline
			{\bf Total} & {\bf 109}
		\end{tabular}
		&
 	   	\begin{tabular}{R{2cm}|L{0.8cm}} 
			Verbal Reasoning & 160+\\
			Quantitative Reasoning & 170\\
			Analytical Writing & 4+
		\end{tabular}
	\end{tabular}
\end{table}
{\renewcommand{\arraystretch}{1}


{\color[HTML]{444444}
\textsuperscript{$*$} {\it Only the paper-based GRE test is administered in Iran (results: \textbf{Dec 9\textsuperscript{th}}). This table is my \uline{expected} scores based on a few (\href{http://store.ets.org/store/ets/en_US/pd/ThemeID.12805600/productID.253480200}{official} and \href{http://www.ets.org/gre/revised_general/prepare/powerprep2/}{ETS--provided}) practice tests.}
\\[8pt]
\textsuperscript{$**$} {\it I'm taking another iBT test (results: \textbf{Nov 12\textsuperscript{th}}) to improve my score in the Speaking section.}
}











\needspace{0.5\textheight}%
    \newdimen\boxwidth%
    \boxwidth=\dimexpr\linewidth-5\fboxsep\relax%
    \colorbox[HTML]{BEDBE3}{%
    \begin{tabularx}{\boxwidth}{c|X}
	    \Writinghand & 
	    {365, E2, Ekbatan, Tehran, Iran}\smallskip\\
	    \Telefon     & 
	    {\plus{ }98 (935) 431 26 45}\smallskip\\
	    \Letter      & 
	    \href{mailto:pooriaazimi@gmail.com}{pooriaazimi@gmail.com}
	    \\[3pt]
%	    \Mundus
	     & 
%
%		\href{https://github.com/pooriaazimi/}{\includegraphics[width=13pt]{Logos/GitHub.pdf}}
		\href{https://github.com/pooriaazimi/}{\includegraphics{Logos/GitHub.eps}}
%
%		\href{http://stackoverflow.com/users/347353/pooria-azimi}{\includegraphics[width=12pt]{Logos/StackOverflow.pdf}}
		\href{http://stackoverflow.com/users/347353/pooria-azimi}{\includegraphics{Logos/StackOverflow.eps}}
%
%		\href{https://news.ycombinator.com/user?id=pooriaazimi}{\includegraphics[width=12pt]{Logos/HN.pdf}}
		\href{https://news.ycombinator.com/user?id=pooriaazimi}{\includegraphics{Logos/HN.eps}}
%		
		\href{http://forum.irmug.com}{\includegraphics[width=14pt]{Logos/IRMUG.pdf}}
%		
	\end{tabularx}}
	





\vspace{8pt}











\section{Teaching Experiences}

\begin{ta}

\item{Database Design}{Winter '11}

\item{Database Design}{Winter '12}

\item{Artificial Intelligence}{Winter '12}

\item{Database Design}{Winter '13}

\item{Data Mining}{Spring '14}

\end{ta}



\vspace{15pt}


\begin{presenter}

\item{Spring '12}
	 {4th National Linux Festival}
	 {\href{http://bit.ly/linux2012}{Advanced Unix Command-line}}

\item{Spring '12}
	 {AUT-CEIT Computing Festival}
	 {Introduction to Node.js}


\item{Spring '13}
	 {AUT Database Workshop}
	 {\href{http://bit.ly/aut-postgresql}{PostgreSQL vs. MySQL}}

\item{Spring '13}
	 {AUT Database Workshop}
	 {\href{http://bit.ly/aut-mongodb}{Introduction to MongoDB}}

\item{Spring '13}
	 {AUT Database Workshop}
	 {Basics of Neo4j and Redis}
\end{presenter}


\vspace{-15pt}











\section{Technical skills}

\hspace{-1pt}%
\textsf{\textbf{Programming Languages}}
	 
\hspace{2pt}%
Experienced:
\mbox{\href{http://www.oracle.com/technetwork/java/}{\bf Java} -- %
     \href{https://www.ruby-lang.org/en/}{\bf Ruby} -- %
     \href{https://en.wikipedia.org/wiki/Objective-C}{\bf Objective-C} -- %
     \href{http://coffeescript.org}{\bf CoffeeScript}%
}

\hspace{2pt}%
Familiar:
\mbox{\href{http://www.scala-lang.org}{\bf Scala} -- %
     \href{https://en.wikipedia.org/wiki/C_(programming_language)}{\bf C} -- %
     \href{http://www.erlang.org}{\bf Erlang} \& %
     \href{http://elixir-lang.org}{\bf Elixir} --  %
     \href{http://schemers.org}{\bf Scheme} -- %
     \href{http://php.net}{\bf PHP}%
}



\vspace{9pt}%
\hspace{-1pt}%
\textsf{\textbf{Databases}}
	 
\hspace{2pt}%
Used extensively: %
     \href{http://www.mongodb.org}{\bf MongoDB} -- %
     \href{http://www.postgresql.org}{\bf PostgreSQL} -- %
     \href{http://redis.io}{\bf Redis}%

\hspace{2pt}%
Used occasionally: %
     \href{http://www.neo4j.org}{\bf Neo4j} -- %
     \href{http://www.mysql.com}{\bf MySQL} -- %
     \href{http://www.microsoft.com/en-us/sqlserver/default.aspx}{\bf MS SQL Server}%



\vspace{9pt}%
\hspace{-1pt}%
\textsf{\textbf{Server-side Web development}}
 
\hspace{2pt}%
Experienced: %
\mbox{\href{http://nodejs.org}{\bf Node.js} -- %
     \href{http://rubyonrails.org}{\bf Ruby on Rails} -- %
     \href{http://www.sinatrarb.com}{\bf Sinatra} -- %
     \href{http://php.net}{\bf PHP}%
}



\vspace{9pt}%
\hspace{-1pt}%
\textsf{\textbf{Client-side Web development}}
	 
\hspace{2pt}%
Limited experience: %
     \href{http://backbonejs.org}{\bf Backbone.js} -- %
     \href{http://emberjs.com}{\bf Ember.js}%



\vspace{9pt}%
\hspace{-1pt}%
\textsf{\textbf{Miscellaneous}}

\hspace{2pt}%
Source Control Management (%
	\href{http://git-scm.com}{\bf git} \& %
	\href{http://www.selenic.com/mercurial/}{\bf hg}%
)%

\hspace{2pt}%
UNIXish tools (%
	\href{http://www.vim.org}{\bf vi}, %
	\href{http://awk.info}{\bf awk}, %
	\href{http://www.gnu.org/software/sed/}{\bf sed}, etc.%
)%

\hspace{2pt}%	
	\href{http://nginx.org/en/}{\bf nginx} -- %
	\href{http://www.webkit.org}{\bf WebKit} -- %
	\href{http://www.antlr.org}{\bf ANTLR} -- %
	\href{http://en.wikipedia.org/wiki/Behavior-driven_development}{\bf BDD} -- %
	\href{http://www.latex-project.org}{\bf \LaTeX}




\newpage





\section{Notable Academic Projects}

\begin{project_list}

\item{July 2011 -- April 2012}
     {Iran Telecommunications Research Center}
     {Kavandeh Search Engine}
	 {Improving link-based Web page ranking algorithms in a Persian-only search engine, using various statistical and heuristic methods.
	 \\[1.5pt]
	 Our team of two heavily improved upon Nutch's \mbox{scoring}, parsing, crawling, and spam detection submodules, and used WebKit's rendering engine for detecting \mbox{semantically} significant parts of a (rendered) Web page and \mbox{assigning} more weight to links in such areas (this feature proved to be computationally expensive and was eventually dropped).}
     {%
     	Apache \href{http://nutch.apache.org}{Nutch} \&
		\href{http://lucene.apache.org/solr/}{Solr} -- %
		\href{http://www.webkit.org}{WebKit} -- %
		\href{http://www.oracle.com/technetwork/java/}{Java} -- %
		\href{https://en.wikipedia.org/wiki/C\%2B\%2B}{C\plus\plus}%
	}

\vspace{-6pt}

\item{March 2012 -- June 2012}
     {}
     {Visual WebPage Segmentation {\large (VPS)}}
	 {Detecting Web page structure with statistical \mbox{analysis} of the visual representation of the rendered page \mbox{content}, and using that structure for improving ranking algorithms in a search engine.\\[1.5pt]%
\mbox{\small\it(development halted after 4 months due to time constraints)}%
}
	 {%
	 	\href{http://nodejs.org}{Node.js} -- %
		\href{http://phantomjs.org}{PhantomJS} -- %
		\href{http://www.mongodb.org}{MongoDB}%
	 }

\vspace{-6pt}

\item{January 2011 -- June 2011}
	 {}
	 {\href{https://github.com/baygan/Baygan}{Baygan Database}}
	 {An extendable and clearly-modulated framework for introducing students to the intricacies of relational database design. Inspired by \href{http://www.stanford.edu/class/cs140/projects/pintos/pintos.html}{pintos}.\\[1.5pt]%
\mbox{\small\it(development halted after 6 months due to time constraints)}%
}
	 {%
	 	\href{http://www.oracle.com/technetwork/java/}{Java} -- %
		\href{http://www.antlr.org}{ANTLR}%
	 }

\vspace{-6pt}
	 
\item{June and July 2010}
	 {}
	 {Embedded Search Engine}
	 {A complete, single-purpose search engine (all parts \mbox{written} from scratch), designed to use minimal online memory (60 MB), for indexing and searching the contents of the English Wikipedia.}
	 {%
		 \href{http://www.oracle.com/technetwork/java/}{Java}%
	 }

\end{project_list}



\vspace{-10pt}










\section{References}

\begin{description}
\item[Shahram Khadivi] (Assistant Professor, Amirkabir \\University of Technology, Iran)
\\
\href{http://ceit.aut.ac.ir/~khadivi/}{http://ceit.aut.ac.ir/{\tilde}khadivi/} -- \href{mailto:khadivi@aut.ac.ir}{Email}
\end{description}
%
\vspace{-3pt}
For much of the past 3 years, and also the coming year, I have been working closely with Dr.\ Khadivi on \mbox{multiple} projects (Kavandeh, VPS, Baygan), Berlin Data \mbox{Mining} Cup, and now my thesis project. I have also been Teacher Assistant to a couple of Dr.\ Khadivi's courses (``Database Design'': three years, ``Data Mining'': one semester).
\\[6pt]
\uline{More references are available upon request} -- but I've \mbox{included} only Dr.\ Khadivi's address above because I think he would be in the best position to evaluate my academic and technical expertise.








\section{Other Projects}

\begin{project_list}
	 
\item{2010 -- 2013}
	 {}
	 {Open Source Contributions}
	 {Contributing to multiple Open Source projects (code, documentation, bug report, and IRC support), \mbox{including} %
	 \href{http://phantomjs.org}{PhantomJS}, %
	 \href{https://github.com/ggreer/the_silver_searcher}{Ag (the silver searcher)}, %
	 \href{http://bitbucket.org/sjl/hg-prompt/}{hg-prompt}, %
	 \href{https://github.com/allending/Kiwi}{Kiwi}, %
	 \href{https://github.com/clowwindy/shadowsocks/}{shadowsocks}, %
	 \href{http://fishshell.com}{fish shell}, and %
	 \href{https://github.com/mxcl/homebrew}{Homebrew}.}%
	 {%
	 	\href{https://developer.mozilla.org/en-US/docs/Web/JavaScript}{JavaScript} -- %
	 	\href{http://awk.info}{awk} -- %
		\href{http://python.org}{Python} -- %
		\href{https://en.wikipedia.org/wiki/Objective-C}{Objective-C} -- %
		\href{https://www.ruby-lang.org/en/}{Ruby}%
	 }
	 
\end{project_list}

\begin{oneline_projects}

\item{\href{http://pooriaazimi.github.io/BetterDictionary/}{BetterDictionary}}
	 {%
	 	\href{https://en.wikipedia.org/wiki/Objective-C}{Objective-C} -- %
		\href{https://en.wikipedia.org/wiki/Cocoa_(API)}{Cocoa}%
	 }	

\item{\href{https://github.com/pooriaazimi/Mini-Java}{MiniJava Parser}}
	 {%
		 \href{http://coffeescript.org}{CoffeeScript} -- %
		 \href{http://zaach.github.io/jison/}{jison} -- %
		 \href{http://d3js.org}{d3.js}%
	 }

\item{\href{https://github.com/pooriaazimi/twitter}{A Twitter Clone!}}
	 {%
		 \href{http://rubyonrails.org}{Ruby on Rails} -- %
		 \href{http://www.postgresql.org}{PostgreSQL} -- %
		 \href{http://coffeescript.org}{CoffeeScript} -- %
		 \href{http://sass-lang.com}{Sass}%
	 }


\item{\href{http://www.turnedondigital.com/?portfolio=farhang-iphone-app}{Farhang} \large{\rm(partial contribution)}}
	 {%
	 	\href{https://en.wikipedia.org/wiki/Objective-C}{Objective-C} -- %
		\href{https://en.wikipedia.org/wiki/Cocoa_Touch}{Cocoa Touch} -- %
		\href{https://en.wikipedia.org/wiki/Core_Data}{Core Data}%
	 }

\item{\href{https://github.com/pooriaazimi/secure_file_system}{Secure File System}}
	 {%
		 \href{http://php.net}{PHP} -- %
		 \href{http://www.mongodb.org}{MongoDB}%
	 }

\item{\href{https://github.com/pooriaazimi/adserver}{Ad Server}}
	 {%
	 	\href{https://www.ruby-lang.org/en/}{Ruby} -- %
		\href{http://www.sinatrarb.com}{Sinatra} -- %
		\href{http://www.sqlite.org}{SQLite}%
	 }


\end{oneline_projects}





\vspace{-20pt}




\section{Online Education}

In addition to attending my normal classes, I have watched the videos, and finished the assignments, of a dozen freely-available online courses, \mbox{including} the following CS-related ones:

\begin{list}{\labelitemi}{\leftmargin=1em}

\item MIT's \href{http://groups.csail.mit.edu/mac/classes/6.001/abelson-sussman-lectures/}{\mbox{Structure and Interpretation of Computer Programs}} {\it(1986 -- by \href{https://en.wikipedia.org/wiki/Hal_Abelson}{Harold Abelson} and \href{https://en.wikipedia.org/wiki/Gerald_Jay_Sussman}{Gerald Jay Sussman})}

\item UC Berkley's \href{http://www.cs.berkeley.edu/~kubitron/courses/cs162-F08/}{Operating Systems and Systems Programming} {\it(2008)}

\item Harvard's \href{http://cs50.tv/2010/fall/}{Introduction to Computer Science} {\it(2010)} and \href{http://cs75.tv/2010/fall/}{Building Dynamic Websites} {\it(2010)}

\item Stanford's \href{http://see.stanford.edu/see/courseinfo.aspx?coll=824a47e1-135f-4508-a5aa-866adcae1111}{Programming Methodology} {\it(2007)}, \href{http://see.stanford.edu/see/courseinfo.aspx?coll=11f4f422-5670-4b4c-889c-008262e09e4e}{\mbox{Programming} Abstractions} {\it(2008)}, \href{http://see.stanford.edu/see/courseinfo.aspx?coll=2d712634-2bf1-4b55-9a3a-ca9d470755ee}{\mbox{Programming} Paradigms} {\it(2008)}, and \href{http://www.stanford.edu/class/cs193p}{iPhone Application Programming} {\it(2011 \& 2013)}
	
\end{list}

I'm also taking the following Coursera courses this semester: 

\begin{list}{\labelitemi}{\leftmargin=1em}

\item EPFL's \href{https://www.coursera.org/course/progfun}{Functional Programming Principles in Scala} {\it(2013 -- by \href{https://en.wikipedia.org/wiki/Martin_Odersky}{Martin Odersky})}

\item EPFL/Typesafe Inc.'s \href{https://www.coursera.org/course/reactive}{Principles of Reactive\\ Programming} {\it(2013 -- by \href{https://en.wikipedia.org/wiki/Martin_Odersky}{Martin Odersky}, et al.)}

\end{list}








\end{document}
